\documentclass[11pt,a4paper]{article}

\usepackage[margin=2.5cm]{geometry}
\usepackage[T1]{fontenc}
\usepackage{xcolor}
\usepackage{listings}
\usepackage{upquote}
\usepackage[hidelinks]{hyperref}
\usepackage{booktabs}

\definecolor{sqlkeyword}{RGB}{0,70,160}
\definecolor{sqlstring}{RGB}{148,0,110}
\definecolor{sqlcomment}{RGB}{0,120,0}
\definecolor{sqlbg}{RGB}{244,244,244}
\definecolor{sqlframe}{RGB}{190,190,190}

% Grey box, coloured tokens. Keywords are written in uppercase, matching
% the SQL that the repair and expansion prompts ask the model to emit.
\lstdefinestyle{sql}{
  language=SQL,
  basicstyle=\ttfamily\small,
  keywordstyle=\color{sqlkeyword}\bfseries,
  ndkeywordstyle=\color{sqlkeyword}\bfseries,
  stringstyle=\color{sqlstring},
  commentstyle=\color{sqlcomment}\itshape,
  backgroundcolor=\color{sqlbg},
  frame=single,
  rulecolor=\color{sqlframe},
  framesep=6pt,
  xleftmargin=2pt,
  xrightmargin=2pt,
  breaklines=true,
  breakatwhitespace=false,
  showstringspaces=false,
  columns=fullflexible,
  keepspaces=true,
  tabsize=2,
  upquote=true,
  sensitive=false,
  captionpos=b,
  aboveskip=10pt,
  belowskip=8pt,
  morekeywords={VARCHAR,INTEGER,PRIMARY,KEY,AUTO_INCREMENT,LIMIT,OFFSET,
    COUNT,AVG,MAX,MIN,SUM,GROUP,ORDER,JOIN,LEFT,RIGHT,INNER,OUTER,
    DISTINCT,EXCEPT,INTERSECT,UNION,BETWEEN,LIKE,ASC,DESC}
}

\lstset{style=sql}

\title{Upgrading \texttt{sql-create-context-copy}:\\
Validation, Repair, Schema Expansion, and Hardness}
\author{Thomas Bayer\\
RWU Hochschule Ravensburg-Weingarten\\
\texttt{thomas.bayer@rwu.de}
\and
Grok}
\date{Working paper, October 2026}

\begin{document}
\maketitle

\begin{abstract}
This note documents a database-free upgrade of the 78{,}577-row text-to-SQL
dataset \texttt{sql-create-context-copy}~\cite{sqlcontextcopy}. String literals
are rewritten to single quotes, every row is checked for SQL syntax, DDL
syntax, and schema consistency, 541 failing rows are repaired, and the DDL of
the valid rows is expanded without changing the query. Each row then receives
an extended hardness label. Where the official Spider adapter accepts the
statement, it also receives a Spider hardness label. No statement is executed.
\end{abstract}

The code and the prompts are in \url{https://github.com/bayerth/sql_tools}.
The refined rows, including both hardness fields, are in
\path{data/sql/sql_context_v2.json}.

Repair and DDL expansion call \texttt{openai/gpt-oss-120b} through LM Studio on
a Mac Studio with an M4 Max and 128\,GB of memory. The context length is
32{,}768 tokens and the temperature is 0. Expansion sends 16 rows per request.
Repair sends one row per request.

\section{The source dataset}

\texttt{sql-create-context-copy}~\cite{sqlcontextcopy} is the local starting
point, saved from the Hugging Face dataset
\texttt{philschmid/sql-create-context-copy} at revision \texttt{ed48b9fa}.
Each row has a natural-language \texttt{question}, a SQL \texttt{answer}, and
a \texttt{context} of one or more \texttt{CREATE TABLE} statements. The copy
contains 78{,}577 training rows. The published schemas are small: they usually
contain only the columns the query mentions, and string columns are often
declared as bare \texttt{VARCHAR}.

The upgrade keeps the questions and, except for the repaired rows, the SQL
answers. It changes the form of literals, the consistency of identifiers, the
richness of the DDL, and the labels attached to each row. Structural
validation is not a check that the SQL answers the question. Nothing in this
pipeline opens a database connection.

\section{Pipeline}

The scripts were run in this order. Each step reads the previous step's output.

\begin{enumerate}
  \item Rewrite double-quoted string literals in \texttt{answer} to single quotes.
  \item Validate SQL syntax, DDL syntax, and schema references with the generic
        SQLGlot dialect.
  \item Repair the 541 rows that failed that validation. A candidate is kept
        only when the repaired SQL and DDL pass validation again.
  \item Merge the 78{,}036 rows that were already valid with the 541 repaired
        rows.
  \item Expand the DDL. The model may add columns and lengths. It must not
        change the SQL.
  \item Merge expanded, unchanged, and failed expansion rows into
        \path{data/sql/sql_context_v2.json}, 78{,}577 rows. Failed rows keep
        their original DDL.
  \item Assign Spider hardness and extended hardness on that file.
\end{enumerate}

\section{Single-quoted literals}
\label{sec:quotes}

The copy stores many string values in double quotes. On MySQL without
\texttt{ANSI\_QUOTES}, and on SQLite, those quotes are string literals. The
generic SQLGlot dialect reads them as identifiers, so a later schema check
reports a missing column. That false error would have dominated validation.

The module \texttt{rewrite\_sql\_string\_quotes} rewrites paired double quotes
in \texttt{answer} only.
\texttt{question} and \texttt{context} are copied unchanged. An apostrophe
inside a literal is doubled, which is the SQL escape for a single quote. A
double-quoted name written directly after \texttt{FROM} is left alone: three
rows use \texttt{FROM "t"} as a table identifier rather than a string. The
default input is \path{data/sql_context_copy}. The default output is
\path{data/sql/sql_context_copy_single_quoted.json}.

On this copy the rewrite changed 62{,}967 statements and doubled an apostrophe
in 503 places. The three quoted \texttt{FROM} names stayed quoted.

\begin{lstlisting}[caption={Before: a double-quoted literal is an identifier in the generic dialect.}]
SELECT Census_Ranking
FROM city
WHERE Status <> "Village"
\end{lstlisting}

\begin{lstlisting}[caption={After: the same predicate, with a single-quoted string.}]
SELECT Census_Ranking
FROM city
WHERE Status <> 'Village'
\end{lstlisting}

\begin{lstlisting}[caption={Before: an apostrophe inside a double-quoted literal.}]
SELECT candidates
FROM table_1342393_17
WHERE incumbent = "James O'Connor"
\end{lstlisting}

\begin{lstlisting}[caption={After: the apostrophe is doubled inside a single-quoted literal.}]
SELECT candidates
FROM table_1342393_17
WHERE incumbent = 'James O''Connor'
\end{lstlisting}

\section{Validation}

Validation uses SQLGlot and does not execute SQL. A row is valid only when
three checks pass: the query parses, the DDL parses, and every name in the
query resolves against that DDL. The generic dialect is the right dialect for
this pass. The copied DDL uses bare \texttt{VARCHAR}. A MySQL check requires
an explicit length and then skips schema resolution, so a MySQL run would
reject almost every row before the interesting errors appear.

The generic pass accepted 78{,}036 rows and rejected 541. Among the rejected
rows, 187 have invalid SQL, 191 have valid syntax but a missing or ambiguous
column, and 163 have a schema that cannot be checked, typically because the
DDL repeats a table or column definition.

\subsection{DML syntax}

A numeric prefix, the keyword \texttt{AS}, and an identifier suffix inside an
aggregate are not a legal argument. The generic parser stops at \texttt{AS}.
The DDL of this row is fine, so schema resolution is not attempted.

\begin{lstlisting}[caption={Invalid SQL. The aggregate argument is not an expression.}]
SELECT COUNT(1 AS st_prize__) AS $__
FROM table_11621747_1
WHERE tournament = 'The Transamerica'
\end{lstlisting}

\begin{lstlisting}[caption={DDL for the same row. Generic syntax accepts bare VARCHAR.}]
CREATE TABLE table_11621747_1 (tournament VARCHAR)
\end{lstlisting}

\subsection{DDL syntax}

Bare \texttt{VARCHAR} is legal in the generic dialect and illegal in MySQL,
which requires a length. The same department query passes the generic DDL
check and fails the MySQL DDL check. Schema resolution does not run after
that failure. This is why the upgrade validates with the generic dialect
first and treats the MySQL length as a later normalisation, not as a reason
to drop the row.

\begin{lstlisting}[caption={Query that is syntactically valid in both dialects.}]
SELECT name, born_state, age
FROM head
ORDER BY age
\end{lstlisting}

\begin{lstlisting}[caption={DDL rejected by MySQL: VARCHAR needs a length.}]
CREATE TABLE head (
  name VARCHAR,
  born_state VARCHAR,
  age VARCHAR
)
\end{lstlisting}

\subsection{Schema references}

Here the query and the DDL both parse. The filter column is absent from every
table, so the schema check fails. Validation reports the missing name. It does
not invent a column.

\begin{lstlisting}[caption={Valid SQL that names a column the DDL does not define.}]
SELECT t1.catalog_name, t1.date_of_publication
FROM catalogs AS t1
JOIN catalog_structure AS t2
  ON t1.catalog_id = t2.catalog_id
WHERE catalog_level_number > 5
\end{lstlisting}

\begin{lstlisting}[caption={The DDL has no catalog\_level\_number column.}]
CREATE TABLE catalogs (
  catalog_name VARCHAR,
  date_of_publication VARCHAR,
  catalog_id VARCHAR
);
CREATE TABLE catalog_structure (
  catalog_id VARCHAR
)
\end{lstlisting}

\section{Repair}

Only the 541 invalid rows are sent to the model. Rows that already validated
stay out of this step. The target dialect of the repair prompt is MySQL. The
model may reconstruct a corrupted identifier, quote a reserved word, merge
column names that differ only in letter case, and give bare \texttt{VARCHAR}
the length 255 when it rewrites a table. A candidate is applied only when the
new SQL and the new DDL pass validation. Passing that check means the
statement is structurally consistent. It does not prove that the repair
preserved the meaning of the question.

The prize query from the DML example is a corrupted column written as an
aggregate argument. The repair turns the suffix into a quoted alias and leaves
the tournament filter in place. The count expression becomes \texttt{COUNT(1)}.
That is a legal statement. It is not evidence that the original table had a
column of that name; the source table is not available.

\begin{lstlisting}[caption={Repaired prize query. The alias is quoted, and the aggregate parses.}]
SELECT COUNT(1) AS `st_prize__`
FROM table_11621747_1
WHERE tournament = 'The Transamerica'
\end{lstlisting}

The catalog query fails because \texttt{catalog\_level\_number} is missing.
The repair keeps the query, quotes the filter column, adds it to
\texttt{catalog\_structure}, and writes explicit string lengths.

\begin{lstlisting}[caption={Repaired catalog DDL. The filter column is now defined.}]
CREATE TABLE catalogs (
  catalog_name VARCHAR(255),
  date_of_publication VARCHAR(255),
  catalog_id VARCHAR(255)
);
CREATE TABLE catalog_structure (
  catalog_id VARCHAR(255),
  `catalog_level_number` INTEGER
)
\end{lstlisting}

The merged file after this step has 78{,}036 originally valid rows plus 541
repaired rows.

\section{Expansion}

Expansion runs only after the statement is valid. The model may add plausible
columns and companions that the published schema omitted, such as \texttt{long}
next to \texttt{lat}, and it gives bare \texttt{VARCHAR} an explicit length
when it rewrites a table. The SQL query must stay byte-for-byte the same. A
changed DDL is kept only when the original SQL still resolves against it. If
the model leaves the DDL unchanged, or if validation rejects the candidate,
the original pair is kept.

Of the 78{,}577 rows, 68{,}983 received a new DDL, 8{,}025 were left unchanged,
and 1{,}569 failed. The failed rows stay in the dataset with their previous
DDL. Added columns are synthetic. They make the schema less of a giveaway of
the query, and they are not columns from a production database.

\begin{lstlisting}[caption={SQL. Expansion is required to leave this statement unchanged.}]
SELECT name
FROM station
WHERE lat < 37.5
\end{lstlisting}

\begin{lstlisting}[caption={DDL before expansion. Only the columns used by the query are present.}]
CREATE TABLE station (
  name VARCHAR,
  lat INTEGER
)
\end{lstlisting}

\begin{lstlisting}[caption={DDL after expansion. A companion column and string lengths were added.}]
CREATE TABLE station (
  name VARCHAR(255),
  lat INTEGER,
  long INTEGER
)
\end{lstlisting}

\section{Run statistics}

The figures below are reconstructed from the request logs of the two model
runs, not from a separate profiler. Each log line records prompt tokens,
completion tokens, reasoning tokens, output tokens, and the request time.
Output tokens are completion tokens plus reasoning tokens. Tokens per second
are output tokens divided by the sum of those request times. Wall-clock time
and seconds per row include everything between the start and the end of the
run. A short expansion probe of 68 requests is not included.

\begin{table}[ht]
  \centering
  \caption{Repair and expansion on the Mac Studio (M4 Max, 128\,GB),
  model \texttt{openai/gpt-oss-120b}, context length 32{,}768.}
  \label{tab:runs}
  \begin{tabular}{lrr}
    \toprule
     & Repair & Expansion \\
    \midrule
    Rows              &       541 &     78{,}577 \\
    Requests          &       541 &      4{,}912 \\
    Rows per request  &         1 &           16 \\
    Wall clock (s)    &   3{,}054.6 & 74{,}452.3 \\
    Seconds per row   &      5.65 &       0.95 \\
    Request time (s)  &   3{,}049.6 & 74{,}293.8 \\
    Prompt tokens     & 1{,}176{,}635 & 12{,}872{,}757 \\
    Completion tokens &   158{,}924 &  3{,}103{,}812 \\
    Reasoning tokens  &    76{,}062 &  2{,}253{,}562 \\
    Output tokens     &   234{,}986 &  5{,}357{,}374 \\
    Output tokens/s   &      77.1 &       72.1 \\
    \bottomrule
  \end{tabular}
\end{table}

The last expansion request contains the remaining row, so the 4{,}912 requests
are 4{,}911 batches of 16 and one batch of 1. Request time is almost the whole
wall clock: about 0.95 seconds per row either way. Repair is slower per row
because each of the 541 statements is its own request.

\section{Hardness}

Hardness is assigned last, on the merged SQL and DDL. Classification does not
call a model and does not execute SQL. \texttt{spider\_hardness} follows the
original Spider rules~\cite{yu2018spider}: \texttt{easy}, \texttt{medium},
\texttt{hard}, and \texttt{extra}. The counters look only at the outermost
adapted statement. A join inside a derived table does not change those
counters.

\texttt{extended\_hardness} is a structural label of this project. It does not
replace the Spider label. A nested \texttt{SELECT} raises the label to at
least \texttt{hard}. It becomes \texttt{extreme} when a nested \texttt{SELECT}
has depth two or more, when at least two nested \texttt{SELECT}s exist and one
of them aggregates, when one nested \texttt{SELECT} is both correlated and
aggregated, or when a nested \texttt{SELECT} is itself Spider \texttt{extra}.
The two labels may be compared only when both are set.

If the official Spider adapter rejects a parsed statement,
\texttt{spider\_hardness} stays null and \texttt{extended\_hardness} is still
set. An expression Spider cannot store counts as one ordinary column, and a
clause it cannot store, such as \texttt{OFFSET}, is ignored. Nesting rules
then apply to that score. A statement that does not parse in the generic
dialect is tried once as MySQL, so backtick identifiers still receive a label.

No row is Spider \texttt{easy} and extended \texttt{extreme}. Six rows are
Spider \texttt{easy} and extended \texttt{hard}. Each of them is a
\texttt{COUNT(*)} over a subquery. Spider scores only the outer statement, and
that statement has no \texttt{WHERE}, grouping, ordering, limit, or join of
its own, so the counters stay at zero and the label is \texttt{easy}. The
extension sees the nested \texttt{SELECT} and therefore assigns at least
\texttt{hard}. One subquery is not enough for \texttt{extreme}.

\begin{lstlisting}[caption={Spider easy, extended hard. The outer count hides a grouped subquery.}]
SELECT COUNT(*)
FROM (
  SELECT city
  FROM airports
  GROUP BY city
  HAVING COUNT(*) > 3
)
\end{lstlisting}

The salary query below is Spider \texttt{extra} because the outer statement
has two nested queries. It is \texttt{extreme} for a narrower reason: two
nested \texttt{SELECT}s, and one of them aggregates.

\begin{lstlisting}[caption={Spider extra, extended extreme. The inner average is one of two nested queries.}]
SELECT employee_id, first_name, last_name, salary
FROM employees
WHERE salary > (SELECT AVG(salary) FROM employees)
  AND department_id IN (
    SELECT department_id
    FROM employees
    WHERE first_name LIKE '%J%'
  )
\end{lstlisting}

\begin{table}[ht]
  \centering
  \caption{Hardness of all 78{,}577 rows after the generic classification run.
  Spider shares use the 75{,}382 rows that received a Spider label. Extended
  shares use every row.}
  \label{tab:hardness}
  \begin{tabular}{lrr}
    \toprule
    Label & Spider & Extended \\
    \midrule
    easy    & 51{,}082 & 53{,}001 \\
    medium  & 22{,}693 & 23{,}957 \\
    hard    &    913   &    921 \\
    extra   &    694   &    690 \\
    extreme &      --- &      8 \\
    \midrule
    labeled & 75{,}382 & 78{,}577 \\
    \bottomrule
  \end{tabular}
\end{table}

The 3{,}195 rows with only an extended label are statements the official
Spider adapter rejects. Eight rows are \texttt{extreme}. No row was left
without an extended label.

Those labels are columns of \path{data/sql/sql_context_v2.json}.

Every row has an \texttt{extended\_hardness} value.
A \texttt{spider\_hardness} value is set on the 75{,}382 rows the official
adapter accepts, and it is null on the other 3{,}195.
\texttt{dialect} is \texttt{generic} on the 78{,}036 rows that passed the
first validation, and \texttt{mysql} on the 541 repaired rows.

\texttt{classification\_status} is \texttt{classified} when both hardness
fields are set, and \texttt{extended} when only the extended label is set.
The last entry of \texttt{status\_history} records that same value.

\section{What the upgrade does not claim}

The questions are the questions of \texttt{sql-create-context-copy}. For
77{,}000-odd rows the SQL answer is unchanged apart from quote normalisation.
The 541 repairs are the rows that did not parse or did not resolve; they are
consistent with their DDL after repair, and their meaning was not checked
against a database. Expanded columns were proposed by a model. Hardness
measures statement shape, not how difficult a person would find the question.

\begin{thebibliography}{9}

\bibitem{sqlcontextcopy}
Philipp Schmid.
\newblock \emph{sql-create-context-copy}.
\newblock Hugging Face dataset, 2023, revision \texttt{ed48b9fa}.
\newblock \url{https://huggingface.co/datasets/philschmid/sql-create-context-copy}.

\bibitem{yu2018spider}
Tao Yu, Rui Zhang, Kai Yang, Michihiro Yasunaga, Dongxu Wang, Zifan Li,
  James Ma, Irene Li, Qingning Yao, Shanelle Roman, Zilin Zhang, and Dragomir Radev.
\newblock Spider: A large-scale human-labeled dataset for complex and
  cross-domain semantic parsing and text-to-SQL task.
\newblock In \emph{Proceedings of the 2018 Conference on Empirical Methods in
  Natural Language Processing}, pages 3911--3921, Brussels, 2018.

\end{thebibliography}

\end{document}
